\section{Builds}
\label{sec:builds}

\paragraph{One source, one compiler, one profile.} The crate is compiled by
a pinned nightly Rust (\code{nightly-2026-10-03}, rustc 1.101.0) with
\code{lto = "fat"}, one codegen unit and \code{panic = "abort"}, for two
targets: x86-64 with its baseline instruction set (SSE2; no
\code{target-cpu=native} anywhere) and \code{wasm32} with the target
features \code{+atomics}, \code{+bulk-memory}, \code{+mutable-globals} and
\code{+simd128}. WASM threads need the standard library rebuilt with
atomics (\code{-Z build-std}), which exists only on nightly; a fixed date
makes the compiler, the library and LLVM the same for both targets. The
crates are locked: Rayon 1.12.0, rayon-core 1.13.0, wasm-bindgen 0.2.129,
wasm-bindgen-rayon 1.3.0. No \code{wasm-opt} pass rewrites the module.

\paragraph{128-bit SIMD on both sides, verified on the artifacts.} WASM has
128-bit SIMD only, so the native build is held to SSE2, also 128-bit. The
code has no explicit SIMD: the vector instructions come from the compiler's
auto-vectorization, so parity is a property to test, not to assume. The
filter's interior loop is a separate, never-inlined function
(\code{blur\_interior}) so that it keeps a symbol in the final binaries, and
the continuous integration disassembles the final native binary and the
final WASM module on every change and requires packed 16-bit arithmetic in
that function on both. With whole-program optimization the vectorization
happens at the final link, so checking a library's assembly would mislead.

\paragraph{What differs below the shared source.} ``The same code'' is
claimed for the crate, not for the layers below it, which are the
environment under study:
\begin{itemize}[leftmargin=1.5em]
\item \emph{Flags.} The target features and link arguments exist only for
WASM; a native run records empty extra flags.
\item \emph{Standard library.} Native links the prebuilt library; WASM
rebuilds it under the project's profile.
\item \emph{Locks inside Rayon.} \code{wasm-bindgen-rayon} enables a Rayon
feature that replaces its mutex and condition variable with spin-based ones,
on WASM only.
\item \emph{Idle threads.} On \code{wasm32} with atomics the standard
library's \code{yield\_now} does nothing, so the yield rounds of Rayon's
idle loop spin without giving up the CPU where native code calls
\code{sched\_yield}; to sleep, a WASM worker waits in
\code{memory.atomic.wait32} and a native worker on a futex.
\item \emph{A second compiler.} The engines compile the WASM for the host
CPU, so they may use encodings newer than the SSE2 the native build is held
to. The vector width is fixed at 128 bits, not the instruction set
(Section~\ref{sec:engine-code}).
\end{itemize}
The idle and wake path is inside every timed iteration of a parallel
region. It is the leading candidate for the browsers' larger fork-join cost
(Section~\ref{sec:forkjoin}); no experiment here isolates it.

\paragraph{The build is reproducible and made twice.} On the prepared
machine the build stage runs the tests and then builds twice from scratch;
the tag \code{measured/v3} (commit \code{ac65af6a}) was made only because
both builds gave the same SHA-256 digests for the native runner
(\code{525707de...}), the WASM module (\code{38787129...}) and the
JavaScript glue. Every campaign records the digests of the binaries it ran,
and refuses to resume with another commit, artifact or browser version.

\subsection{What changed in v3, and what the compiler made of it}
\label{sec:codegen}

A change in the measured path is allowed only with a record of why, and is
checked by a script that builds the base and the new source at the same
path and compares the normalized instructions of every hot function. Before
any change, the source of this repository gave the machine code of v2:
native 88 of 88 hot functions identical, WASM 87 of 87.

\begin{enumerate}[leftmargin=2em]
\item \emph{Every shared-counter histogram builds its tasks through the
grain control}, as the local histogram and the filter do. With Rayon's
adaptive split this calls the library's default, and a test counts the
tasks with and without the call on a pool of one thread: the counts are
equal. The script reports the \code{run} of the three histograms as
changed (it now reads the grain) and the Rayon functions instantiated for
them as renamed. The counting loops were therefore compared by hand, from
just before the first atomic add to just after the last: natively the loop
of 37 instructions is identical (\code{movzbl; shl \$7; lock incq} eight
times per turn, the same registers), and in WASM the loop of 127
instructions is identical. The code that counts is unchanged; the code that
splits now is the same for every histogram.
\item \emph{Sequential references.} Two workloads with no Rayon call, at
one thread only: one loop over the rows that calls \code{blur\_interior},
and one loop that counts into one plain table. Adding a workload can change
another one: the first draft called the parallel filter's row function,
which had one caller and was inlined; with a second caller the compiler
kept it as a call and the machine code of the parallel filter changed,
though its source had not. The sequential filter therefore has its own copy
of the row and edge handling and shares only \code{blur\_interior}; a test
keeps its output equal to the parallel filter's. With both changes, the
parallel filter's task function, \code{blur\_interior}, the local histogram
and the checksum functions are identical to v2 in every variant.
\item \emph{The memory size before the measured phase is written.} It was
already sampled, for the check that fails a WASM run whose memory grows
while it is measured; no new sample is taken. A native filter run shows
what it is for: 69.7~MB before the warm-up, 136.8~MB before the first
measured iteration, 136.8~MB after the last. The 64~MiB of the output
buffer are first touched in the warm-up.
\end{enumerate}

The golden checksums are unchanged. Because the split code changed, the
times of v3 are never compared number by number with those of v2; the
exploratory twin campaigns still show, as a consistency observation, that
the cells whose task code changed kept their times (the private histogram
within 0.2\% of v2 at eight threads, Section~\ref{sec:extras}).
